\documentclass{article}
\usepackage{mathesis}
\usepackage[utf8]{inputenc}
\usepackage[russian]{babel}

\begin{document}
% defined-in: zorich-2
% entity-type: definition
% entity-id: analytic_function

\begin{definition}[Аналитическая функция]
Функция $f: D \to \mathbb{R}$ (где $D \subset \mathbb{R}$) называется аналитической в точке $x_0 \in D$, если существует \entityref{neighborhood}{окрестность} $U(x_0)$ этой точки, в которой функция $f$ может быть представлена сходящимся к ней \entityref{taylor_series}{рядом Тейлора}:
\[
\mForall x \in U(x_0): f(x) = \sum_{k=0}^{\infty} \frac{f^{(k)}(x_0)}{k!} (x - x_0)^k
\]
\end{definition}


% defined-in: zorich-1
% entity-type: definition
% entity-id: taylor_series

\begin{definition}[Ряд Тейлора]
Пусть функция $f$ имеет производные всех порядков в точке $x_0$. Степенной ряд вида
\[
\sum_{k=0}^{\infty} \frac{f^{(k)}(x_0)}{k!} (x - x_0)^k
\]
называется рядом Тейлора функции $f$ в точке $x_0$.
\end{definition}


% defined-in: zorich-1
% entity-type: definition
% entity-id: neighborhood

\begin{definition}[Окрестность точки]
Окрестностью $U(x_0)$ точки $x_0 \in \mathbb{R}$ называется любой интервал $(a, b)$, содержащий точку $x_0$. 
Часто рассматривают $\varepsilon$-окрестность:
\[
U_\varepsilon(x_0) = \{x \in \mathbb{R} \mid |x - x_0| < \varepsilon\}
\]
\end{definition}



\end{document}
